\documentclass[11pt]{article}
\usepackage[T1]{fontenc}
\usepackage{lmodern}
\usepackage[margin=1.08in]{geometry}
\usepackage{amsmath,amssymb,mathtools,microtype}
\usepackage{xurl}
\usepackage[colorlinks=true,linkcolor=blue,urlcolor=blue]{hyperref}
\hypersetup{pdftitle={The irrationality exponent of log(3 + 2 sqrt(2))/sqrt(2) is 2},pdfauthor={Ryan Matthew Casper}}
\setlength{\emergencystretch}{3em}
\setcounter{secnumdepth}{0}
\providecommand{\tightlist}{\setlength{\itemsep}{0pt}\setlength{\parskip}{0pt}}
\title{The irrationality exponent of\par $\log(3+2\sqrt{2})/\sqrt{2}$ is $2$}
\author{Ryan Matthew Casper}
\date{October 7, 2026}
\begin{document}
\maketitle
\textbf{Verification.} All 14 new modules, including the exact endpoint,
passed one coherent Lean 4.34.1 build and a foundational-axiom audit.
Nanoda and a separate stock Lean 4.34.0 replay checked the same frozen
export of 117,830 declarations, with only the three standard
foundational axioms.

\section{1. Statement and
normalization}\label{statement-and-normalization}

Let \[
u=3+2\sqrt2,\qquad x=\frac{\log u}{\sqrt2}.
\] The target is \[
\mu(x)=2.
\tag{1}
\] The proof establishes the stronger quantified statement \[
\forall\nu>2\ \exists Q\in\mathbb Z,\ Q\ge2:\quad
\forall p,q\in\mathbb Z,\ q\ge Q\Longrightarrow
q^{-\nu}\le\left|\frac{\log(3+2\sqrt2)}{\sqrt2}-\frac pq\right|.
\tag{2}
\] Fractions need not be reduced, and the threshold is uniform in the
numerator and denominator. The threshold is not asserted to be
effective. The usual rational irrationality exponent is the supremum of
exponents admitting infinitely many reduced rational approximants with
positive error smaller than the corresponding negative power of the
denominator.

The normalization in (1) is essential to this argument. We prove (2)
directly. We do not infer it by multiplying a theorem about a logarithm
by an irrational algebraic number, and we do not infer the exponent of
the unnormalized logarithm from (1).

Put \(\beta=\sqrt2\), \(\mathcal O=\mathbb Z[\beta]\), and let \[
\sigma_+(a+b\beta)=a+b\sqrt2,\qquad
\sigma_-(a+b\beta)=a-b\sqrt2.
\] These are injective ring homomorphisms into
\(\mathbb R\subset\mathbb C\). The norm is \[
N(a+b\beta)=a^2-2b^2=\sigma_+(a+b\beta)\sigma_-(a+b\beta)\in\mathbb Z.
\] For a nonzero order element its norm is a nonzero integer, so \[
1\le |\sigma_+(z)|\,|\sigma_-(z)|.
\tag{3}
\] Unlike the imaginary-quadratic case, (3) does not bound either
embedding separately.

Here \(N(u)=1\), \(u>1\), and \(\sigma_-(u)=3-2\sqrt2=u^{-1}>0\).
Writing \(\omega=\log u=\beta x\), the two compatible exponentials are
\[
e^{\beta x}=u,\qquad e^{-\beta x}=u^{-1}.
\tag{4}
\] Both sequences of nonnegative powers are injective. The formal
generic input allows \(U\in\mathcal O\), a positive integer \(D\), a
nonzero real \(x\), and the two identities \[
e^{\pm\beta x}=\sigma_\pm(U)/D
\tag{5}
\] with injective powers. These are explicit arithmetic conditions. The
endpoint uses \(U=3+2\beta\), \(D=1\), and verifies every condition.
This package is confined to \(\mathbb Q(\sqrt2)\); it does not certify
every real quadratic field.

\section{2. One selected minor and its two
embeddings}\label{one-selected-minor-and-its-two-embeddings}

Fix \(\nu>2\). Suppose rational approximations \(r_k=p_k/q_k\)
satisfying \[
|x-r_k|\le q_k^{-\nu}
\tag{6}
\] exist with arbitrarily large positive denominators. Choose finitely
many of them, with successively separated weights \[
w_k=\lceil\log q_k\rceil,\qquad v_k=w_k/\theta,\qquad
T_k=\lceil F_0w_k/v_0\rceil,\qquad
G_k(t)=\sum_{1\le l<T_k}\frac{(-1)^{l+1}}l t^l.
\] The positive rational parameters satisfy \(0<\theta<1\) and
\(F_0>2/\theta\). Their order of selection is given in §5.

At sign \(\varepsilon\in\{+1,-1\}\), write
\(\alpha_\varepsilon=e^{\varepsilon\beta x}=\sigma_\varepsilon(U)/D\).
Use multiplicative-additive centers \[
(Y,X_1,\ldots,X_m)=
(\alpha_\varepsilon^j,\varepsilon\beta jr_1,\ldots,\varepsilon\beta jr_m),
\qquad 0\le j<K.
\] Columns are monomials \(Y^hX^\gamma\) with
\(w_0h+\sum_kw_k\gamma_k\le H\). Rows are \((j,s,a)\) with \[
0\le j<K,\qquad v_0s+\theta^{-1}\sum_kw_ka_k<H.
\] The actual matrix entry is \[
A^\varepsilon_{(j,s,a),(h,\gamma)}=
[t^su^a]\,(\alpha_\varepsilon^j(1+t))^h
\prod_k(\varepsilon\beta jr_k+G_k(t)+u_k)^{\gamma_k}.
\tag{7}
\] The factor \(\alpha_\varepsilon^{jh}\) remains part of the matrix.

The inherited logarithmic interpolation theorem applies to arbitrary
complex additive centers and distinct nonzero multiplicative centers.
Its strict weighted-volume, coordinate-ratio, and separated-product
conditions are supplied by §5. It gives surjectivity of \(A^+\) onto all
rows for sufficiently large heights in a cofinal sequence \(H=nR\),
where \(R\) clears the rational weights. Truncation preserves the
relevant jet packets because \(T_kv_0\ge F_0w_k>v_k\). The theorem is
proved by the inherited weighted curve comparison, logarithmic
differential frame, ordinary blowup, and Serre vanishing argument; no
interpolation conclusion is supplied as an endpoint hypothesis.

Select a square nonzero minor \(\Delta_H^+\) containing every row.
Define \(\Delta_H^-\) by using \textbf{exactly the same selected
columns} in \(A^-\). No second minor is selected. Put \[
M_H=\#\{\text{rows}\},\qquad
b_H=\frac{\sum_{\rm rows}\sum_kw_ka_k}{M_HH}.
\] Then \(0\le b_H\le\theta\), and lattice counting gives \[
M_H\sim\frac{K\theta^mH^{m+1}}{(m+1)!v_0\prod_kw_k}.
\tag{8}
\]

\section{3. A common cleared order
matrix}\label{a-common-cleared-order-matrix}

Let \[
L_k=\operatorname{lcm}(1,\ldots,T_k),\qquad
E_H=\prod_kL_k^{\lfloor H/w_k\rfloor}.
\] Multiply each selected column \((h,\gamma)\) by
\(D^{Kh}\prod_kq_k^{\gamma_k}\), each row by \(\prod_kq_k^{-a_k}\), and
every entry by \(E_H\). If \(a\le\gamma\), the resulting order-valued
entry before evaluation is \[
E_H U^{jh}D^{(K-j)h}\binom\gamma a
[t^s](1+t)^h
\prod_k(\beta jp_k+q_kG_k(t))^{\gamma_k-a_k}.
\tag{9}
\] It is zero otherwise. The exponent \((K-j)h\) is nonnegative, and
each \(\gamma_k-a_k\le\lfloor H/w_k\rfloor\), so the \(L_k\) factors
clear all logarithm coefficients. Thus (9) belongs to \(\mathcal O\).

For additional precision, the formal proof first constructs one
polynomial \(Q_T\in\mathbb Z[t]\) with \(Q_T=L_TG_T\). For exponent
\(d\le e\), it expresses the cleared factor as \[
L_T^{e-d}(L_T\beta jp+qQ_T(t))^d.
\] This fixes one order polynomial, coefficient, and matrix before
either embedding is applied. Ring homomorphisms commute with its finite
products, powers, coefficients, and determinant. Evaluation at
\(-\beta\) changes \(U\) and the additive centers while leaving \(Q_T\),
row indices, and selected columns unchanged. Therefore the two evaluated
determinants are precisely \[
\sigma_\pm(\det B_H)=C_H\Delta_H^\pm,
\quad
C_H=E_H^{M_H}
\prod_{\rm columns}\left(D^{Kh}\prod_kq_k^{\gamma_k}\right)
\prod_{\rm rows}\prod_kq_k^{-a_k}>0.
\tag{10}
\] This is an identity for the actual selected determinants. It uses
only the two explicit embeddings of the quadratic order, not an assumed
automorphism of all real or complex numbers.

Since \(\Delta_H^+\ne0\), also \(\det B_H\ne0\). Equation (3) then
proves both \(\Delta_H^-\ne0\) and \[
1\le C_H^2|\Delta_H^+|\,|\Delta_H^-|.
\tag{11}
\] Let \(\Lambda=4+\log4\) and \(w_* =\min_kw_k\). From
\(\log L_k\le\Lambda T_k\), the column degree constraint, and
\(w_k-1\le\log q_k\le w_k\), the normalized arithmetic consequence is \[
\frac{\log|\Delta_H^+|+\log|\Delta_H^-|}{2M_HH}
\ge -(1-b_H)-E_{\rm ar},
\tag{12}
\] where \[
E_{\rm ar}=\frac{\Lambda F_0m}{v_0}
+\Lambda\sum_k\frac1{w_k}+\frac\theta{w_*}
+\frac{K\log D}{w_0}.
\tag{13}
\] Columns cost at most \(M_HH(1+K\log D/w_0)\); row factors save at
least \(M_HH(b_H-\theta/w_*)\). The two embeddings double both the
clearing logarithm and the analytic logarithms, so averaging retains the
original coefficient \(1-b_H\). At our integral-unit endpoint \(D=1\),
the base-denominator term is zero.

\section{4. Analytic smallness at both
embeddings}\label{analytic-smallness-at-both-embeddings}

Choose \(\rho=100\max(1,2|x|)\), which covers both
\(\omega_\varepsilon=\varepsilon\beta x\). For a monomial column define
\[
f_{b,P}(z)=[u^b]P(e^z,z+u_1,\ldots,z+u_m)
=\binom\gamma b e^{hz}z^{|\gamma|-|b|}.
\] For \(|t|<1\), put
\(z_j^\varepsilon(t)=j\omega_\varepsilon+\log(1+t)\). Then
\(e^{z_j^\varepsilon(t)}=\alpha_\varepsilon^j(1+t)\), and the additive
arguments in (7) are \[
z_j^\varepsilon(t)+u_k+e_{jk}^\varepsilon+\tau_k(t),\qquad
e_{jk}^\varepsilon=\varepsilon\beta j(r_k-x),\quad
\tau_k=G_k-\log(1+t).
\tag{14}
\] The rational truncation stays fixed at the negative embedding; its
local comparison is \(-j\omega+\log(1+t)\). No principal-logarithm
identity for powers is needed.

Both signs have the same error norm. Using \(|\beta|\le2\), the finite
row translation expansion gives scalar size at most \[
\exp\{-\nu w(b-a)+HE_{\rm tr}\},\qquad
E_{\rm tr}=\frac\nu{F_0}+\frac{\log2}{v_0}
+\frac{\log4+\log(2K)+\nu+\log2}{w_*}.
\tag{15}
\] Here a nonzero tail term has multiplicities with
\(\sum_ke_kT_k\le s\), hence \(\sum_kw_ke_k<H/F_0\). The additional
slope cost is the single term \(\log2/w_*\), which will be paid
explicitly in §5.

On \(|z|\le\rho K\), the column functions have size bounded by
\(\exp\{H(\rho K/w_0+\log(2\rho K)/w_*)\}\). All comparison points on
\(|t|=1/2\) lie in the disc of radius \(\rho K/2\). Taylor expansion has
an absolutely convergent geometric majorant. When rows with the same
transverse index use the same Taylor order, the determinant term
vanishes. If there are \(n_b\) rows of transverse index \(b\), surviving
orders sum to at least \(\binom{n_b}2\).

With \(c=\log2/4\), the inherited collision estimate supplies the saving
\(-c\sum_bn_b^2\) and normalized holomorphic cost \[
E_{\rm hol}=\frac{\rho K}{w_0}+\frac{\log2}{v_0}
+\frac{\log(2\rho K)}{w_*}.
\] Let \(N_A(H)\) count transverse indices of weight at most \(AH\), and
set \[
c_H=\frac{c\eta^2M_H}{HN_A(H)},\qquad
c_\infty=\frac{c\eta^2K\theta^m}{(m+1)v_0A^m}.
\] If at least \(\eta M_H\) rows use low indices, Cauchy--Schwarz gives
saving \(-c_H\). Otherwise the approximation errors give saving
\(-\nu(A(1-\eta)-b_H)\). The finite row-expansion count is bounded by
\((\lfloor H\rfloor+1)^{2m}(\lfloor H/v_0\rfloor+1)\); its normalized
logarithm and the collision remainder tend to zero. Therefore, uniformly
in the selected columns and at each sign, \[
\frac{\log|\Delta_H^\pm|}{M_HH}
\le E_{\rm an}+o(1)+
\max\{-c_H,-\nu(A(1-\eta)-b_H)\},
\quad E_{\rm an}=E_{\rm tr}+E_{\rm hol},
\tag{16}
\] and \(c_H\to c_\infty\). All weights, approximants, and parameters
are fixed before this height limit. Averaging the two inequalities in
(16) gives the identical upper bound for the left side of (12).

\section{5. Parameters and the
contradiction}\label{parameters-and-the-contradiction}

Replace the lcm budget by \(\Lambda'=\Lambda+\log2\). For every
\(\nu>2\), choose positive rational \(\theta,A,B,C\) such that \[
0<\theta<A<B<1,\quad \nu(A-\theta)>1-\theta,
\quad C>1,\quad CB<1,\quad B<C\theta<1.
\tag{17}
\] For example choose rational \(b\in(1/2,1-1/\nu)\), take
\(\theta=1-s\), \(A=1-bs\) for small positive rational \(s\), then
choose \(C\in(A/\theta,1/A)\) and \(B\in(A,\min(1/C,C\theta))\). The
inequality \(A^2<\theta\) ensures these intervals are nonempty.

Choose small rational \(\eta>0\) with
\(g=\nu(A(1-\eta)-\theta)-(1-\theta)>0\), and let
\(\epsilon=\min(g/2,1/2)\). Choose \(F_0>2/\theta\) with
\(\nu/F_0<\epsilon/3\). As \(m\) grows, set \[
K=\lfloor C^m\rfloor,\qquad w_0=B^{-m},\qquad
v_0=2K\theta^mw_0.
\] Then \(K/w_0\le(CB)^m\to0\), \(v_0\to\infty\), \(m/v_0\to0\), and \[
c_\infty=\frac{c\eta^2(B/A)^m}{2(m+1)}\longrightarrow\infty.
\] Fix a sufficiently large \(m\) so that \[
\frac{\Lambda'F_0m+2\log2}{v_0}
+\frac{(\rho+\log D)K}{w_0}<\epsilon/3,\qquad c_\infty>2.
\tag{18}
\] The interpolation volume is \(K(w_0/v_0)\theta^m=1/2\), while
\(K\theta^m<1\). A small positive rational margin \(\sigma\) gives the
required strict volume and coordinate-ratio conditions.

Only now choose the approximants from (6), with all weights sufficiently
large and successively separated to meet the inherited finite product
comparisons. For equally sized index subsets, common factors cancel;
increasing the next weight supplies the remaining strict comparison. The
weights may also satisfy \[
\Lambda'\sum_k\frac1{w_k}
+\frac{\theta+\log4+\log(2K)+\nu+\log(2\rho K)}{w_*}<\epsilon/3.
\tag{19}
\] The interpolation threshold is chosen after these fixed centers and
weights, so there is no circular dependence on height.

The finite nonempty set of weights attains its minimum, giving the exact
reserve \[
\frac{\log2}{w_*}\le\log2\sum_k\frac1{w_k}
\le(\Lambda'-\Lambda)
\left(\frac{F_0m}{v_0}+\sum_k\frac1{w_k}\right).
\tag{20}
\] Thus (18), (19), and the initial choice of \(F_0\) imply \[
E_{\rm ar}+E_{\rm an}<\epsilon<g,
\qquad c_\infty>1+E_{\rm ar}+E_{\rm an}.
\tag{21}
\]

For large cofinal heights, (12) and the average of (16) are
incompatible. The collision alternative would require
\(c_H\le1-b_H+E_{\rm ar}+E_{\rm an}+o(1)\), contradicting (21). The
approximation alternative would require \[
\nu(A(1-\eta)-b_H)-(1-b_H)
\le E_{\rm ar}+E_{\rm an}+o(1).
\] Its left side is at least \(g\) since \(b_H\le\theta\) and \(\nu>1\),
again contradicting (21). Hence arbitrarily large denominators
satisfying (6) do not exist, proving (2).

If \(x\) were rational, its unreduced exact fractions with unbounded
denominators would violate (2). Thus \(x\) is irrational. Dirichlet
approximation gives exponent at least 2; (2) excludes every larger
exponent, proving (1).

\section{6. Relation to prior work and
provenance}\label{relation-to-prior-work-and-provenance}

The normalized family is established in approximation theory. Bashmakova
and Zolotukhina's \href{https://www.mathnet.ru/eng/cheb531}{2017 primary
article} treats \(\sqrt d\log((\sqrt d+1)/(\sqrt d-1))\); at \(d=2\)
this is exactly twice our target. Their normalization is therefore
directly comparable by rational scaling. Related quadratic logarithms
also appear in \href{https://arxiv.org/abs/1501.06752}{Polyanskii's
work}. We do not identify the strongest previously published numerical
bound for this exact target or claim exhaustive priority.

The proof reuses OpenAI's weighted interpolation framework and released
library at
\href{https://github.com/openai/math/tree/adc7f1241b42e322a6451854ab7e4b4c146bf78a}{commit
adc7f1241b42e322a6451854ab7e4b4c146bf78a}, Mathlib, and the author's
preceding logarithm and imaginary-quadratic adaptations. The new
mechanism is the single order-valued cleared matrix evaluated at both
real embeddings, together with analytic savings for its same selected
minor at both signs. The averaging in (12) removes the degree-two loss
without discarding the row savings. The interpolation geometry, Taylor
collision estimate, parameter existence, and final ordinary-exponent
consequence are substantial inherited results.

\textbf{AI use disclosure.} Ryan Matthew Casper initiated and directed
the investigation. OpenAI coding agents contributed substantially to the
research assessment, mathematical adaptation, Lean proof development,
verification tooling, and exposition. Formal checking concerns the
specified statements relative to their foundational axioms. It is not
independent human peer review. No external referee was contacted for
this work.

\section{Verification materials}\label{verification-materials}

The \href{lean/RealNorm/Main.lean}{Lean endpoint},
\href{lean/VERIFICATION.md}{verification summary}, and
\href{VERIFY.md}{reproduction instructions} accompany this paper. The
proof sources are frozen; verification status and exact export
identities are recorded with the package.
\end{document}
